\begin{scanpage}{0018}{1}
\chaptertitle{I}{DIFFRACTION BY A WEDGE}
As was already said in the Introduction, the field scattered by a body may be investigated in the form of the sum of the fields being radiated by the uniform and nonuniform parts of the surface current. The uniform part of the current is completely determined by the geometry of the body and the magnetic field of the incident wave. The nonuniform part generally is unknown. However, one may approximately assume that in the vicinity of the discontinuity of a convex surface it will be the same as on a corresponding wedge. Therefore, it is necessary for us to begin by studying the diffraction of a plane electromagnetic wave by a wedge. This chapter will be devoted to this problem. First we will investigate the rigorous solution of this problem (\S\ 1 and 2). Then we will find its solution in the physical optics approach (\S\ 3). The difference of these solutions determines the field created by the nonuniform part of the current (\S\ 4).
\sectiontitle{1}{The Rigorous Solution}
The rigorous solution to the problem of diffraction of a plane wave by a wedge was first obtained by Sommerfeld by the method of branching wave functions [16]. Later, the diffraction of cylindrical and spherical waves by a wedge also was studied. A rather extensive bibliography on these problems may be found, for example, in the
\end{scanpage}
